\section{Additional Experimental Results}
\label{appdx:experimental-results}

In this appendix, we provide a small synthetic study validating the intrinsic-dimension scaling behavior predicted by our theory. Our main theoretical claim is that, under the low-dimensional linear latent subspace assumption, the statistically difficult component of the velocity field is the tangent component, whose complexity depends primarily on the intrinsic dimension $d_0$ rather than the ambient dimension $d_x$. The experiments in this appendix test this claim in a controlled setting.

\subsection{Synthetic Setup}

We consider the synthetic latent-subspace model
\begin{align}
h \sim \sum_{m=1}^M \pi_m \mathcal{N}(\mu_m, \Sigma_m),
\qquad
x = Uh,
\end{align}
where $U \in \mathbb{R}^{d_x \times d_0}$ has orthonormal columns and $h \in \mathbb{R}^{d_0}$ follows a Gaussian mixture model with $M=4$ components. We generate $U$ by orthogonalizing a Gaussian random matrix and fix one canonical GMM for each value of $d_0$.

For the forward interpolation path, we use the standard linear flow-matching path
\begin{align}
X_t = t X_1 + (1-t) X_0,
\qquad
X_0 \sim \mathcal{N}(0, I_{d_x}),
\end{align}
with $t \sim \mathrm{Unif}[\varepsilon, 1-\varepsilon]$ and $\varepsilon = 0.01$.

Our latent model receives the projected latent coordinate
\begin{align}
\bar{h} = U^\top x_t \in \mathbb{R}^{d_0}
\end{align}
together with a sinusoidal time embedding and predicts the tangent velocity in latent space. The analytically known normal component is added back exactly:
\begin{align}
\hat{u}(x_t,t)
=
U \hat{v}(\bar{h}, t)
-
\frac{1}{1-t}(I-UU^\top)x_t.
\end{align}

The model architecture is fixed across all settings: a four-layer MLP with hidden width $256$, SiLU activations, AdamW optimizer, learning rate $3 \times 10^{-4}$, batch size $256$, and $150$ epochs. We report averages over five random seeds.

\subsection{Evaluation Metrics}

Our primary metric is the per-dimension tangent oracle MSE:
\begin{align}
\mathrm{MSE}_{\mathrm{tan}}^{\mathrm{per\text{-}dim}}
=
\frac{1}{d_0}
\mathbb{E}
\left[
\left\|
\hat{v}(\bar{h}, t) - v^\star(\bar{h}, t)
\right\|_2^2
\right],
\end{align}
where the oracle tangent velocity is
\begin{align}
v^\star(\bar{h}, t)
=
\frac{1}{t}\bar{h}
+
\frac{1-t}{t}
\nabla_{\bar{h}} \log p_t^h(\bar{h}).
\end{align}

We divide by $d_0$ so that changes in the metric reflect statistical difficulty rather than the number of output coordinates. We also report the total tangent oracle MSE
\begin{align}
\mathrm{MSE}_{\mathrm{tan}}^{\mathrm{total}}
=
\mathbb{E}
\left[
\left\|
\hat{v}(\bar{h}, t) - v^\star(\bar{h}, t)
\right\|_2^2
\right]
\end{align}
as a secondary metric. We additionally verify that the normal MSE—measuring the discrepancy between the analytically reconstructed normal component and its theoretical expression—is numerically negligible across all runs ($\leq 3 \times 10^{-10}$), confirming that the analytical reconstruction is implemented correctly.

\subsection{Intrinsic versus Ambient Dimension Scaling}

Figure~\ref{fig:appendix_scaling} shows per-dimension tangent oracle MSE as a function of training set size for two experimental panels.

\paragraph{Panel A: varying ambient dimension.}
We fix $d_0 = 4$ and vary $d_x \in \{32, 128, 512\}$. To make the comparison as clean as possible, all three values of $d_x$ use the same latent GMM, the same latent training samples, and the same held-out evaluation set; only the ambient embedding matrix $U$ changes across conditions.

The three curves are nearly identical at all training set sizes. At $n_{\mathrm{train}} = 2000$ the spread across $d_x$ values is negligible (ratio of largest to smallest mean MSE: $1.01\times$), and at $n_{\mathrm{train}} = 50000$ it remains small ($1.29\times$), with all three curves lying in the same order of magnitude and overlapping within error bars. There is no systematic trend with $d_x$: the curves do not separate or diverge as ambient dimension increases. This confirms that, once the estimator operates only on latent coordinates with oracle access to $U$, changing the ambient embedding dimension has little effect on tangent estimation accuracy.

\paragraph{Panel B: varying intrinsic dimension.}
We fix $d_x = 128$ and vary $d_0 \in \{2, 4, 8, 16\}$, using one independently sampled canonical GMM per value of $d_0$.

The curves separate clearly according to intrinsic dimension. At $n_{\mathrm{train}} = 2000$ the ordering is perfectly monotone with a roughly $6\times$ gap between $d_0 = 2$ and $d_0 = 16$. At larger training set sizes, $d_0 = 2$ remains consistently the easiest and $d_0 = 16$ consistently the hardest across all seeds; the middle values $d_0 = 4$ and $d_0 = 8$ are close to each other at $n_{\mathrm{train}} = 50000$ (means $0.014$ and $0.015$ respectively, within one standard deviation), which we attribute to the different canonical GMMs drawn for each $d_0$ producing problems of similar difficulty at that sample size. All curves decrease with $n_{\mathrm{train}}$, with higher-$d_0$ curves decreasing more slowly.

\begin{figure}[t]
    \centering
    \includegraphics[width=\textwidth]{figures/figure1.pdf}
    \caption{
    Per-dimension tangent oracle MSE as a function of training set size.
    \emph{Left (Panel A):} fixing $d_0 = 4$ and varying $d_x \in \{32,128,512\}$,
    the three curves are nearly identical at all training set sizes, confirming
    that ambient dimension has little effect on tangent estimation accuracy once
    the estimator operates only on latent coordinates.
    \emph{Right (Panel B):} fixing $d_x = 128$ and varying $d_0 \in \{2,4,8,16\}$,
    larger intrinsic dimension produces higher MSE and slower improvement with
    training set size; $d_0 = 2$ is consistently the easiest and $d_0 = 16$
    consistently the hardest.
    Error bars denote standard deviation across five random seeds.
    }
    \label{fig:appendix_scaling}
\end{figure}

\subsection{Tangent Velocity Diagnostics}

Figure~\ref{fig:appendix_tangent_diag} reports three diagnostics for the representative configuration $(d_x, d_0, n_{\mathrm{train}}) = (128, 4, 50000)$.

\paragraph{Panel A: learned versus oracle tangent velocity.}
We plot learned against oracle tangent velocity coordinates for $5000$ held-out points $(\bar{h}, t)$, coloring each point by latent coordinate index. The scatter lies tightly along the diagonal for all four coordinates, with per-coordinate correlations above $0.99$ and per-coordinate MSE ranging from $0.012$ to $0.019$. There is no visible coordinate-specific failure mode.

\paragraph{Panel B: tangent oracle MSE during training.}
Per-dimension tangent oracle MSE decreases from $0.040$ at epoch $10$ to $0.015$ at epoch $150$, a reduction of approximately $2.7\times$. The training loss decreases in parallel. Both quantities improve steadily throughout training with no sign of stagnation or divergence.

\paragraph{Panel C: latent endpoint distribution.}
We generate $2000$ latent samples by integrating the learned ODE $dh/dt = \hat{v}(h,t)$ from Gaussian noise using Euler integration ($100$ steps), and compare against $2000$ draws from the true latent GMM. The per-coordinate means and standard deviations of the generated and reference distributions match closely. Sliced Wasserstein distance (500 projections) between the two distributions is $0.093$. The joint-PCA scatter shows good overlap between the generated and reference clouds.

\begin{figure}[t]
    \centering
    \includegraphics[width=\textwidth]{figures/figure2_dx128_d04_n50000_seed0.pdf}
    \caption{
    Tangent-component diagnostics for the representative setting $(d_x,d_0,n_{\mathrm{train}}) = (128,4,50000)$.
    \emph{Left (Panel A):} learned tangent velocities align closely with oracle tangent velocities across all four latent coordinates, with per-coordinate correlations above $0.99$.
    \emph{Middle (Panel B):} per-dimension tangent oracle MSE decreases $2.7\times$ over $150$ epochs, alongside the flow-matching training loss.
    \emph{Right (Panel C):} generated latent samples match the reference GMM distribution closely; sliced Wasserstein distance is $0.093$.
    }
    \label{fig:appendix_tangent_diag}
\end{figure}

\subsection{Summary}

Panel A shows that once subspace geometry is provided and the estimator operates in latent coordinates, changing $d_x$ over a $16\times$ range produces negligible variation in per-dimension tangent oracle MSE. Panel B shows that increasing $d_0$ makes the problem harder, with a clear separation between $d_0 = 2$ and $d_0 = 16$ at all training set sizes and all curves decreasing with $n_{\mathrm{train}}$. The diagnostic panels confirm that the learned tangent velocity field closely matches the oracle and that the latent ODE generates a distribution consistent with the true GMM. Taken together, these results support the tangent/normal decomposition and the intrinsic-dimension scaling behavior predicted by the theory.
